\section{Conclusion}
\label{sec:conclusion}

This paper asked a question that aggregate evaluation of neuro-symbolic legal reasoning pipelines leaves unanswered: when an LLM extractor feeding a symbolic rule engine makes a predicate error, does that specific error actually change the resulting legal decision? Building on a neuro-symbolic pipeline that couples LLM predicate extraction with the PYTHEN rule tree for hearsay classification, and on an oracle experiment confirming the rule tree is logically sound given correct predicates, we introduced a predicate-level intervention procedure: correcting one erroneous predicate at a time, rerunning the rule tree, and recording whether the decision flips. This yields a Decision Impact Rate (DIR) for every predicate, which we compared directly against standalone extraction F1, and used to refine an extraction-error taxonomy for this pipeline along a rule-aware, decision-relevance dimension.

The results confirm the paper's central hypothesis: extraction difficulty and decision importance are not the same thing, and the rule structure determines how far apart they can be. \textit{is\_written\_assertion} is both the hardest predicate to extract and the one whose errors are most likely to flip the final decision, while \textit{is\_non\_verbal\_assertive\_conduct} is comparably hard to extract but structurally redundant, with errors that rarely reach the output. The refined taxonomy further showed that nominally identical error types carry different practical risk across models - Gemma 4 31B's hallucinated predicates are decision-critical 96.2\% of the time, versus 46.9-59.5\% for the other two models - a distinction invisible to standard predicate-level F1 or accuracy.

Taken together, these findings support treating rule-aware error attribution, rather than aggregate predicate accuracy alone, as a necessary component of evaluating and verifying LLM-symbolic legal reasoning pipelines: knowing that an extractor is imperfect is far less actionable than knowing which of its imperfections the final decision actually depends on. We reported this as a case study on a single LegalBench task with a specific rule structure, and Section~\ref{sec:discussion} outlines how the methodology should be extended to rule trees of different shapes, to joint (multi-predicate) intervention, and to concrete verification tooling built on the Decision Impact Rate. We hope this diagnose-which-errors-matter perspective proves useful beyond the hearsay setting studied here, wherever an LLM's outputs are consumed by an auditable symbolic reasoner rather than trusted directly.
